\documentclass[11pt]{article}
\usepackage[T1]{fontenc}
\usepackage[utf8]{inputenc}
\usepackage{lmodern}
\usepackage[margin=1in]{geometry}
\usepackage{amsmath,amssymb,amsthm,mathtools}
\usepackage{graphicx,booktabs,array}
\usepackage{xcolor}
\usepackage[colorlinks=true,linkcolor=blue!55!black,citecolor=blue!55!black,urlcolor=blue!55!black]{hyperref}
\hypersetup{pdftitle={The 2n-1 steady-state bound is sharp for sequential distributive phosphorylation},
 pdfsubject={Maximal multistationarity of the n-site phosphorylation cycle for every n},
 pdfkeywords={multisite phosphorylation, multistationarity, mass-action kinetics, chemical reaction networks, steady states, interlacing polynomials, formal verification}}

\theoremstyle{plain}
\newtheorem{theorem}{Theorem}[section]
\newtheorem{lemma}[theorem]{Lemma}
\newtheorem{proposition}[theorem]{Proposition}
\newtheorem{corollary}[theorem]{Corollary}
\newtheorem{conjecture}[theorem]{Conjecture}
\theoremstyle{definition}
\newtheorem{definition}[theorem]{Definition}
\newtheorem{example}[theorem]{Example}
\theoremstyle{remark}
\newtheorem{remark}[theorem]{Remark}
\numberwithin{equation}{section}

\newcommand{\R}{\mathbb R}
\newcommand{\Q}{\mathbb Q}
\newcommand{\raw}{\mathrm{raw}}
\newcommand{\ET}{E_{\mathrm T}}
\newcommand{\FT}{F_{\mathrm T}}
\newcommand{\ST}{S_{\mathrm T}}
\newcommand{\JS}{\mathcal J}
\newcommand{\cs}{c_\ast}
\newcommand{\es}{e_\ast}
\DeclareMathOperator{\sign}{sign}
\setlength{\emergencystretch}{1.5em}

\title{The \(2n-1\) steady-state bound is sharp for\\ sequential distributive phosphorylation}
\author{}
\date{}

\begin{document}
\maketitle

\begin{abstract}
The sequential distributive \(n\)-site phosphorylation cycle, in which one kinase and one
phosphatase act on a substrate with \(n\) ordered sites, is the standard mass-action model of
multisite protein modification. Wang and Sontag proved that it has at most \(2n-1\) positive
steady states for any rate constants and any conserved totals. Whether this bound is attained
has been known only for \(n\le 4\). We prove that it is attained for every \(n\). More
precisely, for any \(2n-1\) distinct positive numbers we construct, by explicit algebraic
formulas, positive rate constants and totals for which the system has exactly \(2n-1\) positive
steady states whose free-kinase to free-phosphatase ratios are the prescribed numbers. All of
these steady states are nondegenerate, so the maximal count persists on a nonempty open subset
of the full \((6n+3)\)-dimensional parameter space, and rational data suffice. The construction
rests on a square-root change of variable that turns the steady-state equation on a special
parameter locus into an interpolation problem, on a positive two-term polynomial recurrence, and
on an interlacing argument that gives an explicit admissible range of enzyme totals; for
instance, equally spaced auxiliary roots work for every kinase-to-phosphatase ratio above
\((4n^2-1)/8\). We also prove a determinant formula which shows that \(n-1\) of the constructed
steady states are unstable for every \(n\), and we certify in exact rational arithmetic that
for \(n\le 10\) the remaining \(n\) are asymptotically stable. The existence of \(2n-1\)
positive steady states for every \(n\) is formally verified in Lean~4 with Mathlib.
\end{abstract}

\section{Introduction}

\subsection{Background}
Reversible phosphorylation at multiple sites is one of the most common ways in which cells
regulate protein function \cite{Cohen2000,SalazarHofer2009}. In the simplest mechanistic
description, a single kinase \(E\) adds phosphate groups to a substrate one site at a time in a
fixed order, a single phosphatase \(F\) removes them in the reverse order, and each enzyme
releases the substrate after every catalytic event. This is the \emph{sequential distributive}
mechanism, also called the multiple futile cycle. Under mass-action kinetics it is the reaction
network
\begin{equation}\label{eq:network}
\begin{aligned}
 S_{i-1}+E &\ \underset{b_i}{\overset{a_i}{\rightleftharpoons}}\ C_i
       \ \xrightarrow{\ c_i\ }\ S_i+E,\\[3pt]
 S_i+F &\ \underset{\beta_i}{\overset{\alpha_i}{\rightleftharpoons}}\ Y_i
       \ \xrightarrow{\ \gamma_i\ }\ S_{i-1}+F,
       \qquad i=1,\dots,n,
\end{aligned}
\end{equation}
where \(S_i\) is the substrate with \(i\) phosphorylated sites, \(C_i\) and \(Y_i\) are the
kinase--substrate and phosphatase--substrate complexes, and all \(6n\) rate constants are
positive. There are \(3n+3\) species and three conserved totals,
\begin{equation}\label{eq:totals}
 \ET=E+\sum_{i=1}^n C_i,\qquad
 \FT=F+\sum_{i=1}^n Y_i,\qquad
 \ST=\sum_{i=0}^n S_i+\sum_{i=1}^n(C_i+Y_i).
\end{equation}

Markevich, Hoek, and Kholodenko observed that already the two-site cycle can be bistable, with
no feedback other than competition for the shared enzymes \cite{Markevich2004}. This started a
long line of work on the steady states of \eqref{eq:network}; see
\cite{Gunawardena2005,ConradiFlockerziRaisch2008,ThomsonGunawardena2009,FeliuWiuf2012,
HolsteinFlockerziConradi2013,ConradiMincheva2014,HellRendall2015,CFMW2017,BihanDickensteinGiaroli2020,
FKdWY2022,TorresFeliu2021} and the surveys \cite{ConradiShiu2018,Dickenstein2016,Dickenstein2020}.
The network has become a standard test case for algebraic and analytic methods in reaction
network theory, because it is large, biologically meaningful, and has a steady-state variety
with a monomial parametrization \cite{PDSC2012}.

The basic counting result is due to Wang and Sontag.

\begin{theorem}[Wang and Sontag {\cite{WangSontag2008}}]\label{thm:WS}
For every \(n\ge1\), all positive rate constants, and all positive totals \(\ET,\FT,\ST\), the
system \eqref{eq:network} has at most \(2n-1\) positive steady states with those totals.
\end{theorem}

This is Theorem~3 of \cite{WangSontag2008}. In the same paper they showed that \(n+1\) (for
\(n\) even) or \(n\) (for \(n\) odd) positive steady states do occur, and that, under a
genericity condition, there are at most \(n+1\) when the enzymes are sufficiently scarce
relative to the substrate. On this basis they conjectured that \(n+1\) is never exceeded. For
\(n=2\) the numbers \(n+1\) and \(2n-1\) agree. Flockerzi, Holstein, and Conradi
\cite{FHC2014} refuted the conjecture: they produced rate constants with five steady states
for \(n=3\) and seven for \(n=4\), so that the bound \(2n-1\) is attained for \(n\le4\). Their
method prescribes zeros of a determining equation and requires the solution of a system of
polynomial inequalities that grows with \(n\). Since then the bound \(2n-1\) has been
conjectured to be sharp for all \(n\) \cite{FRW2020,FKdWY2023}, and the question is recorded
as open beyond \(n=4\) in \cite{FKdWY2023}. Thomson and Gunawardena had given numerical
evidence that the number of steady states grows with \(n\) \cite{ThomsonGunawardena2009}.
Feliu, Rendall, and Wiuf \cite{FRW2020} proved that for every \(n\) there are parameters with
\(2\lfloor n/2\rfloor+1\) steady states of which \(\lfloor n/2\rfloor+1\) are asymptotically
stable, and explained why their reduction to a Michaelis--Menten limit cannot give \(2n-1\)
for \(n\ge3\). Giaroli, Rischter, P\'erez Mill\'an, and Dickenstein \cite{GRPD2019} described
explicit open parameter regions with \(2\lfloor n/2\rfloor+1\) positive steady states, and
remarked that degeneration techniques are unlikely to reach \(2n-1\). Thus for every \(n\ge5\)
there has been a gap between the largest number known to occur, \(2\lfloor n/2\rfloor+1\), and
the upper bound \(2n-1\), and the gap grows linearly in \(n\). The dynamics of the cycle
beyond steady states is also of current interest \cite{CMS2019,CFM2020,ConradiMincheva2024}.

\subsection{Results}
This paper closes the gap.

\begin{theorem}\label{thm:main}
Let \(n\ge1\) and let \(u_0,\dots,u_{2n-2}\) be distinct positive real numbers. There exist
positive rate constants \(a_i,b_i,c_i,\alpha_i,\beta_i,\gamma_i\) and positive totals
\(\ET,\FT,\ST\) such that the system \eqref{eq:network} has exactly \(2n-1\) positive steady
states with these totals, and the values of the ratio \(E/F\) of free kinase to free
phosphatase at these steady states are exactly \(u_0,\dots,u_{2n-2}\). Every one of these
steady states is nondegenerate. If the numbers \(\sqrt{1+8u_j}\) are all rational, then the
rate constants, the totals, and all steady-state concentrations can be chosen rational.
\end{theorem}

Here a positive steady state is \emph{nondegenerate} if the Jacobian of the mass-action vector
field, restricted to the stoichiometric subspace, is invertible there. Nondegeneracy makes the
count stable under perturbation, and the upper bound holds for all parameters, so we obtain a
full-dimensional region of maximal multistationarity.

\begin{corollary}\label{cor:open}
For every \(n\ge1\) there is a nonempty open subset of \(\R^{6n+3}_{>0}\), the space of all rate
constants and totals, on which \eqref{eq:network} has exactly \(2n-1\) positive steady states,
all nondegenerate.
\end{corollary}

In particular the maximum number of positive steady states of the \(n\)-site sequential
distributive cycle is \(2n-1\) for every \(n\), and the set of parameters attaining it has
positive Lebesgue measure.

The construction is explicit. Write \(x_j=\sqrt{1+8u_j}\). A two-term recurrence produces two
polynomials \(N,J\) of degree \(n-1\) with positive coefficients from the \(x_j\)
(Section~\ref{sec:recurrence}); an exact polynomial division depending on one more number
\(r\) converts them into the polynomials \(B,D\) whose coefficients are, up to the
normalization \eqref{eq:rates}, the rate constants (Section~\ref{sec:conversion}); and the
totals are \((\ET,\FT,\ST)=(2r,2,2r+2)\). The only condition is that the coefficients come out
positive. We determine an explicit range of \(r\) for which they do.

\begin{theorem}\label{thm:threshold}
In the setting of Theorem~\ref{thm:main}, put \(x_j=\sqrt{1+8u_j}\). The construction of
Sections~\ref{sec:recurrence}--\ref{sec:assembly} yields positive rate constants, and hence
the conclusions of Theorem~\ref{thm:main} with \(\ET/\FT=r\), for every real number \(r\) such
that
\begin{equation}\label{eq:rcondition}
 r>\max_j u_j\qquad\text{and}\qquad 4r+\sqrt{1+8r}\ \ge\ \sum_{j=0}^{2n-2}x_j .
\end{equation}
For \(x_j=j+2\) the second condition follows from the first, so every
\(r>(4n^2-1)/8\) is admissible.
\end{theorem}

Finally we obtain information on stability. Order the prescribed ratios increasingly,
\(u_0<u_1<\dots<u_{2n-2}\).

\begin{theorem}\label{thm:unstable}
In the construction of Theorem~\ref{thm:threshold}, the \(n-1\) steady states with odd index
\(j\) are unstable: the Jacobian on the stoichiometric subspace has an odd number of positive
real eigenvalues there. At the \(n\) steady states with even index this number is even.
\end{theorem}

This follows from a determinant formula, Proposition~\ref{prop:det}, which is valid for all
rate constants and relates the determinant of the restricted Jacobian at any positive steady
state to the slope of a scalar function. For the rational family \(x_j=j+2\) with the smallest
admissible integer \(r\) we have computed the characteristic polynomials of the restricted
Jacobians exactly and applied the Routh criterion in rational arithmetic
(Proposition~\ref{prop:census}): for every \(n\le10\) the \(n\) even-indexed steady states are
asymptotically stable and the \(n-1\) odd-indexed ones are hyperbolic saddles with a
one-dimensional unstable manifold. This is \(n\)-fold multistability, the largest number
compatible with \(2n-1\) alternating steady states; for \(n\ge3\) it exceeds the number
\(\lfloor n/2\rfloor+1\) established for all \(n\) in \cite{FRW2020}. We conjecture that it
holds for every \(n\) (Conjecture~\ref{conj:stable}).

\subsection{Idea of the proof}
At a positive steady state all concentrations are monomials in three quantities: the ratio
\(u=E/F\), the free unphosphorylated substrate \(s=S_0\), and \(f=F\). The three totals then
become \(\ET=uf(1+sB(u))\), \(\FT=f(1+suD(u))\), \(\ST=sA(u)+usf(B(u)+D(u))\) for polynomials
\(A,B,D\) with positive coefficients determined by the rate constants, and conversely every
such triple of polynomials comes from rate constants (Proposition~\ref{prop:param}). The two
enzyme totals determine \(s\) and \(f\) as rational functions of \(u\), and the steady states
are the solutions of one scalar equation \(H(u)=\ST\).
The equation is rational of high degree, and the difficulty is that its \(3n\) effective
coefficients must all be positive while \(2n-1\) roots are placed in the region where \(s\)
and \(f\) are positive.

The key observation is that on the locus \(A=(1+u)D\), \(\ST=\ET+\FT\), the numerator of
\(H(u)-\ST\) is a quadratic form in \(\ell=B/D-r\) and \(v=r-u\) with discriminant
\(v^2(1+8u)\). In the variable \(x=\sqrt{1+8u}\) it splits into two linear factors, and the
steady-state equation becomes the \emph{linear} condition
\begin{equation}\label{eq:introinterp}
 (x-1)\bigl(B(u)-rD(u)\bigr)=2(r-u)D(u),\qquad u=\tfrac{x^2-1}{8}.
\end{equation}
We then need polynomials \(B,D\) in \(u\) of degree \(n-1\) with positive coefficients for
which \eqref{eq:introinterp} holds at \(2n-1\) prescribed values of \(x\). Splitting
\(\prod_j(x-x_j)\) into its even and odd parts in \(x\) gives polynomials in \(u\) that solve a
version of this interpolation problem with \(r=\infty\); they are generated by a recurrence
whose coefficients are visibly positive. A B\'ezout-type quotient moves the solution to finite
\(r\). Positivity of the quotient is where interlacing enters: the roots of \(N\) and \(J\)
are real, negative, and interlace, so \(N/J\) has a partial-fraction expansion with positive
residues, and this expresses \(D\) and \(B\) as positive combinations of polynomials with
positive coefficients.

\subsection{Formal verification}
The existence statement underlying Theorem~\ref{thm:main}, that for every \(n\ge1\) there are
positive rate constants and totals with \(2n-1\) distinct positive steady states of the
mass-action system \eqref{eq:network}, has been formalized and checked in Lean~4
\cite{Lean4} with Mathlib \cite{Mathlib}. The formal statement refers to the vector field
itself and not to a reduced equation. Section~\ref{sec:formal} states precisely what is and is
not covered. All numerical statements in the paper are exact rational computations that are
reproduced by a script distributed with the source.

\subsection{Organization}
Section~\ref{sec:model} parametrizes the steady states and reduces them to a scalar equation.
Section~\ref{sec:recurrence} introduces the square-root coordinate and the recurrence.
Section~\ref{sec:conversion} contains the conversion and the positivity theorem.
Section~\ref{sec:assembly} assembles the steady states and proves the counting part of
Theorem~\ref{thm:main}. Section~\ref{sec:nondeg} proves nondegeneracy, the determinant formula,
Corollary~\ref{cor:open}, and Theorem~\ref{thm:unstable}. Section~\ref{sec:properties} records
further properties of the constructed steady states. Section~\ref{sec:examples} gives a
complete three-site example and the stability computations.
Section~\ref{sec:formal} describes the formal verification, and
Section~\ref{sec:discussion} discusses limitations and open problems.

\section{Steady states of the mechanism}\label{sec:model}
Let \(V\) be the mass-action vector field of \eqref{eq:network} on \(\R^{3n+3}\), with
coordinates
\[
 (S_0,\dots,S_n,\;E,\;F,\;C_1,\dots,C_n,\;Y_1,\dots,Y_n),
\]
and write \(W\colon\R^{3n+3}\to\R^3\) for the linear map \eqref{eq:totals}. The totals are conserved, so \(V\)
takes values in \(\ker W\). We shall see in Section~\ref{sec:nondeg} that \(\ker W\) is the
stoichiometric subspace; it has dimension \(3n\). A \emph{positive steady state} is a zero of
\(V\) with all coordinates positive.

For given rate constants put
\begin{equation}\label{eq:pq}
 p_i=\frac{a_i}{b_i+c_i},\qquad q_i=\frac{\alpha_i}{\beta_i+\gamma_i},\qquad
 \lambda_i=\frac{c_ip_i}{\gamma_iq_i},\qquad t_i=\prod_{k=1}^{i}\lambda_k\quad(t_0=1),
\end{equation}
and define three polynomials with positive coefficients,
\begin{equation}\label{eq:coeff}
\begin{gathered}
 A(u)=\sum_{i=0}^n t_i u^i,\qquad
 B(u)=\sum_{i=0}^{n-1} B_i u^i,\qquad
 D(u)=\sum_{i=0}^{n-1} D_i u^i,\\
 B_{i-1}=p_it_{i-1},\qquad D_{i-1}=q_it_i\qquad(1\le i\le n).
\end{gathered}
\end{equation}
The letter \(D\) always denotes this polynomial; the phosphatase complexes are called \(Y_i\).

\begin{proposition}\label{prop:param}
\textup{(i)} The positive steady states of \eqref{eq:network} are exactly the points
\begin{equation}\label{eq:states}
\begin{aligned}
 S_i&=t_i u^i s \quad(0\le i\le n),& E&=uf,& F&=f,\\
 C_i&=B_{i-1}u^i sf,& Y_i&=D_{i-1}u^i sf \quad(1\le i\le n),
\end{aligned}
\end{equation}
with \((u,s,f)\in\R^3_{>0}\). Their totals are
\begin{equation}\label{eq:totalcoords}
 \ET=uf\bigl(1+sB(u)\bigr),\qquad \FT=f\bigl(1+suD(u)\bigr),\qquad
 \ST=sA(u)+usf\bigl(B(u)+D(u)\bigr).
\end{equation}
\textup{(ii)} Conversely, let \(t_0=1\) and let \(t_1,\dots,t_n\), \(B_0,\dots,B_{n-1}\),
\(D_0,\dots,D_{n-1}\) be arbitrary positive numbers. For any positive numbers
\(b_i,\beta_i,\gamma_i\) \((1\le i\le n)\), the rate constants
\begin{equation}\label{eq:ratesgeneral}
 c_i=\gamma_i\frac{D_{i-1}}{B_{i-1}},\qquad
 a_i=(b_i+c_i)\frac{B_{i-1}}{t_{i-1}},\qquad
 \alpha_i=(\beta_i+\gamma_i)\frac{D_{i-1}}{t_i}
\end{equation}
are positive and have the given numbers as their coefficients \eqref{eq:coeff}.
\end{proposition}

\begin{proof}
The components of \(V\) for the complexes are
\(\dot C_i=a_iS_{i-1}E-(b_i+c_i)C_i\) and \(\dot Y_i=\alpha_iS_iF-(\beta_i+\gamma_i)Y_i\). For
\(1\le i\le n\) let
\begin{equation}\label{eq:pools}
 Z_i=\sum_{j=i}^n S_j+\sum_{j=i+1}^n C_j+\sum_{j=i}^n Y_j
\end{equation}
be the amount of substrate carrying at least \(i\) phosphate groups, where \(C_j\) is counted
with \(j-1\) groups and \(Y_j\) with \(j\). Binding and unbinding do not change \(Z_i\), the
kinase catalytic step at site \(i\) increases it, and the phosphatase catalytic step at site
\(i\) decreases it; hence
\begin{equation}\label{eq:poolcurrent}
 \dot Z_i=c_iC_i-\gamma_iY_i .
\end{equation}
At a steady state, therefore, \(C_i=p_iS_{i-1}E\), \(Y_i=q_iS_iF\), and \(c_iC_i=\gamma_iY_i\),
which gives \(S_i=\lambda_i(E/F)S_{i-1}\). With \(u=E/F\), \(s=S_0\), \(f=F\) this is
\eqref{eq:states}. Conversely, at a point of the form \eqref{eq:states} the complex
components of \(V\) vanish and \(c_iC_i=\gamma_iY_i\) for every \(i\). Since then
\(a_iS_{i-1}E-b_iC_i=c_iC_i\) and \(\alpha_iS_iF-\beta_iY_i=\gamma_iY_i\), we get
\(\dot S_i=-c_{i+1}C_{i+1}+c_iC_i-\gamma_iY_i+\gamma_{i+1}Y_{i+1}=0\) (with the convention
that terms with index \(0\) or \(n+1\) are absent), and \(\dot E=\dot F=0\) because the
complexes are stationary. Summation gives \eqref{eq:totalcoords}. Part (ii) is a direct
substitution: \eqref{eq:ratesgeneral} gives \(p_i=B_{i-1}/t_{i-1}\), \(q_i=D_{i-1}/t_i\), and
\(\lambda_i=c_ip_i/(\gamma_iq_i)=t_i/t_{i-1}\).
\end{proof}

Thus the set of positive steady states, and the number of them in each compatibility class,
depend on the \(6n\) rate constants only through the \(3n\) coefficients of \(A,B,D\). The
remaining \(3n\) degrees of freedom \(b_i,\beta_i,\gamma_i\) affect the dynamics, including
stability, but not the steady states. For definiteness we shall use the normalization
\begin{equation}\label{eq:rates}
 \beta_i=\gamma_i=1,\qquad b_i=c_i=\frac{D_{i-1}}{B_{i-1}},\qquad
 a_i=\frac{2D_{i-1}}{t_{i-1}},\qquad \alpha_i=\frac{2D_{i-1}}{t_i},
\end{equation}
which is \eqref{eq:ratesgeneral} with \(\gamma_i=\beta_i=1\) and \(b_i=c_i\).

\subsection*{The scalar equation}
Fix positive totals and put \(r=\ET/\FT\),
\begin{equation}\label{eq:LM}
 L(u)=B(u)-rD(u),\qquad M(u)=B(u)-uD(u).
\end{equation}
Dividing the first two equations of \eqref{eq:totalcoords} gives \(r-u=suL(u)\). Hence a
positive steady state with \(u\ne r\) satisfies \((r-u)L(u)>0\), and then
\begin{equation}\label{eq:chart}
 s(u)=\frac{r-u}{uL(u)},\qquad
 f(u)=\frac{\FT}{1+s(u)uD(u)}=\frac{\FT L(u)}{M(u)},
\end{equation}
where we used \(L+(r-u)D=M\); in particular \(M(u)=L(u)\bigl(1+suD(u)\bigr)\) is nonzero and
has the sign of \(L(u)\). Conversely, for every \(u>0\) with \((r-u)L(u)>0\) the formulas
\eqref{eq:chart} define positive \(s,f\) such that \eqref{eq:states} has enzyme totals
\(\ET,\FT\). On the open set \(\mathcal U=\{u>0:(r-u)L(u)>0\}\) define
\begin{equation}\label{eq:H}
 H(u)=s(u)A(u)+u\,s(u)f(u)\bigl(B(u)+D(u)\bigr).
\end{equation}
The positive steady states with totals \((\ET,\FT,\ST)\) and \(u\ne r\) correspond bijectively
to the solutions \(u\in\mathcal U\) of \(H(u)=\ST\). This reduction is the one used in
\cite{WangSontag2008}; we need it only as a tool for constructing steady states and for
computing slopes, and we take the upper bound from Theorem~\ref{thm:WS}.

\section{The square-root coordinate and a positive recurrence}\label{sec:recurrence}

\subsection{A locus on which the scalar equation factors}
\begin{lemma}\label{lem:factor}
Suppose \(A(u)=(1+u)D(u)\) and \(\ST=\ET+\FT\), and normalize \(\FT=2\), so that \(\ET=2r\) and
\(\ST=2(r+1)\). For \(u\in\mathcal U\) put \(\ell=B(u)/D(u)-r\), \(v=r-u\), and
\(x=\sqrt{1+8u}\). Then
\begin{equation}\label{eq:residual}
 H(u)-\ST=\frac{(1+u)\bigl(v(\ell+v)-2u\ell^2\bigr)}{u\,\ell\,(\ell+v)}
 =-\frac{2(1+u)}{\ell(\ell+v)}\Bigl(\ell-\frac{2v}{x-1}\Bigr)\Bigl(\ell+\frac{2v}{x+1}\Bigr).
\end{equation}
\end{lemma}
\begin{proof}
Since \(L=D\ell\) and \(M=D(\ell+v)\), \eqref{eq:chart} gives \(s=v/(uD\ell)\) and
\(f=2\ell/(\ell+v)\). Therefore
\[
 H=\frac{(1+u)v}{u\ell}+\frac{2v(\ell+r+1)}{\ell+v},\qquad
 H-2(r+1)=\frac{(1+u)v}{u\ell}-\frac{2\ell(r+1-v)}{\ell+v}
 =(1+u)\Bigl(\frac{v}{u\ell}-\frac{2\ell}{\ell+v}\Bigr),
\]
because \(r+1-v=1+u\). This is the first expression. For the second, expand the product and
use \(x^2-1=8u\):
\(-2u\bigl(\ell^2-\tfrac{4v}{x^2-1}\ell-\tfrac{4v^2}{x^2-1}\bigr)=-2u\ell^2+v\ell+v^2\).
\end{proof}

On \(\mathcal U\cap\{u<r\}\) we have \(\ell,v>0\), so only the first linear factor in
\eqref{eq:residual} can vanish. On this locus the steady-state equation is therefore
equivalent to
\begin{equation}\label{eq:interp}
 \frac{B(u)}{D(u)}=r+\frac{2(r-u)}{x-1},\qquad x=\sqrt{1+8u},
\end{equation}
which is \eqref{eq:introinterp}. The two conditions defining the locus have a simple meaning:
\(A=(1+u)D\) says that \(t_i=D_i+D_{i-1}\), and \(\ST=\ET+\FT\) says that there is as much
substrate as enzyme. They are used only to find a center for the construction;
Corollary~\ref{cor:open} removes both.

\subsection{The recurrence}
Put \(m=n-1\). Let \(x_0,\dots,x_{2m}\) be distinct real numbers greater than \(1\) and write
\begin{equation}\label{eq:ux}
 u(x)=\frac{x^2-1}{8},\qquad u_j=u(x_j),\qquad \Pi(x)=\prod_{j=0}^{2m}(x-x_j).
\end{equation}
Start with \(N=1\) and \(J=(x_0-1)/2\), and for \(k=0,\dots,m-1\), with
\((a,b)=(x_{2k+1},x_{2k+2})\), replace \(N,J\) simultaneously by
\begin{equation}\label{eq:recurrence}
\begin{aligned}
 N_{\mathrm{new}}&=\bigl[8u+(a+1)(b+1)\bigr]N+2(a+b)J,\\
 J_{\mathrm{new}}&=4(a+b)\,uN+\bigl[8u+(a-1)(b-1)\bigr]J .
\end{aligned}
\end{equation}

\begin{lemma}\label{lem:recurrence}
After \(m\) steps \(N=\sum_{i=0}^m h_iu^i\) and \(J=\sum_{i=0}^m d_iu^i\) have degree \(m\),
all coefficients \(h_i,d_i\) \((0\le i\le m)\) are positive, and
\begin{equation}\label{eq:product}
 (x-1)\,N(u(x))-2\,J(u(x))=\Pi(x).
\end{equation}
Moreover
\begin{equation}\label{eq:leading}
 h_m=8^m,\qquad d_m=8^m\,\frac{\sum_jx_j-1}{2},\qquad d_0=J(0)=\frac12\prod_{j}(x_j-1).
\end{equation}
\end{lemma}
\begin{proof}
Let \(N,J\) have positive coefficients in degrees \(0,\dots,k\) and none above. The
coefficient of \(u^i\) in \(N_{\mathrm{new}}\) is \(8h_{i-1}+(a+1)(b+1)h_i+2(a+b)d_i\) and in
\(J_{\mathrm{new}}\) it is \(4(a+b)h_{i-1}+8d_{i-1}+(a-1)(b-1)d_i\), with \(h_{-1}=d_{-1}=
h_{k+1}=d_{k+1}=0\). Since \(a,b>1\), both are positive for \(0\le i\le k+1\). The leading
coefficients obey \(h\mapsto8h\) and \(d\mapsto4(a+b)h+8d\), which gives the first two
formulas in \eqref{eq:leading}. For \eqref{eq:product}, the initial pair gives \(x-x_0\), and a
direct expansion using \(8u=x^2-1\) shows
\[
 (x-1)N_{\mathrm{new}}-2J_{\mathrm{new}}
 =\bigl[x^2-(a+b)x+ab\bigr]\bigl[(x-1)N-2J\bigr]
\]
at \(u=u(x)\): the coefficient of \(N\) is \((x-1)(x^2+ab+a+b)-(a+b)(x^2-1)=(x-1)(x-a)(x-b)\),
and the coefficient of \(J\) is \(2(a+b)(x-1)-2(x^2+ab-a-b)=-2(x-a)(x-b)\). Setting \(x=1\) in
\eqref{eq:product} gives \(d_0\).
\end{proof}

\begin{remark}\label{rem:evenodd}
Write \(\Pi(x)=\Pi_{\mathrm e}(x^2)+x\,\Pi_{\mathrm o}(x^2)\). Replacing \(x\) by \(-x\) in
\eqref{eq:product} and solving the two equations shows that, as functions of \(w=x^2=1+8u\),
\begin{equation}\label{eq:evenodd}
 N=\Pi_{\mathrm o},\qquad J=-\tfrac12\bigl(\Pi_{\mathrm e}+\Pi_{\mathrm o}\bigr).
\end{equation}
So \(N\) and \(J\) are determined by \eqref{eq:product} and depend only on the set
\(\{x_j\}\), not on its ordering. The recurrence is a convenient way to compute them and to
see that their coefficients are positive.
\end{remark}

\section{Conversion and coefficient positivity}\label{sec:conversion}
Let \(N,J\) be as in Lemma~\ref{lem:recurrence}. For \(r>0\) write \(N_r=N(r)\), \(J_r=J(r)\)
and define
\begin{equation}\label{eq:conversion}
 Q(u)=\frac{u\,J_rN(u)-r\,N_rJ(u)}{u-r},\qquad
 D_{\raw}=4Q,\qquad
 B_{\raw}=rD_{\raw}+(4rN_r-2J_r)N-2J_rJ .
\end{equation}
The numerator of \(Q\) vanishes at \(u=r\), so \(Q\) is a polynomial, of degree \(m\) with
leading coefficient \(J_rh_m\).

\begin{lemma}\label{lem:identity}
Let \(\cs=J_r\) and \(\es=4rN_r-J_r\). Then, with \(u=u(x)\),
\begin{equation}\label{eq:interpolation}
 (x-1)\bigl(B_{\raw}(u)-rD_{\raw}(u)\bigr)+2(u-r)D_{\raw}(u)=(\cs x+\es)\,\Pi(x).
\end{equation}
In particular \eqref{eq:interp} holds for \(B_{\raw},D_{\raw}\) at every \(u_j\) with
\(D_{\raw}(u_j)\ne0\).
\end{lemma}
\begin{proof}
We have \(B_{\raw}-rD_{\raw}=(\es-\cs)N-2\cs J\) and
\(2(u-r)D_{\raw}=8\cs uN-2(\cs+\es)J\), because \(8rN_r=2(\cs+\es)\). With \(8u=x^2-1\) the
left side of \eqref{eq:interpolation} becomes
\((x-1)\bigl[\es-\cs+\cs(x+1)\bigr]N-2(\cs x+\es)J=(\cs x+\es)\bigl[(x-1)N-2J\bigr]\), and
\eqref{eq:product} applies.
\end{proof}

It remains to decide when \(D_{\raw}\) and \(B_{\raw}\) have positive coefficients. The answer
comes from the location of the roots of \(N\) and \(J\).

\begin{lemma}[Interlacing]\label{lem:interlace}
Let \(m\ge1\). The polynomials \(N\) and \(J\) have only real, simple roots
\(\nu_m<\dots<\nu_1\) and \(\zeta_m<\dots<\zeta_1\), and
\begin{equation}\label{eq:interlace}
 \nu_m<\zeta_m<\dots<\nu_2<\zeta_2<\nu_1<\zeta_1<0 ,\qquad \nu_1<-\tfrac18 .
\end{equation}
Consequently
\begin{equation}\label{eq:pf}
 \frac{N(u)}{J(u)}=\alpha+\sum_{k=1}^m\frac{\sigma_k}{u-\zeta_k},\qquad
 \alpha=\frac{h_m}{d_m}=\frac{2}{\sum_jx_j-1}>0,\qquad \sigma_k>0 .
\end{equation}
\end{lemma}
\begin{proof}
For \(y>0\) evaluate \eqref{eq:product} at \(x=iy\), where \(u=-(1+y^2)/8\) is real and runs
over \((-\infty,-\tfrac18)\). Each factor is \(iy-x_j=-(x_j^2+y^2)^{1/2}e^{-i\phi_j}\) with
\(\phi_j=\arctan(y/x_j)\). The number of factors is odd, so
\(\Pi(iy)=R\,(-\cos\theta+i\sin\theta)\) with \(R=\prod_j(x_j^2+y^2)^{1/2}>0\) and
\(\theta(y)=\sum_j\phi_j\). Comparing real and imaginary parts in
\((iy-1)N-2J=\Pi(iy)\) gives
\begin{equation}\label{eq:trig}
 N=\frac{R\sin\theta}{y},\qquad J=\frac{R}{2}\Bigl(\cos\theta-\frac{\sin\theta}{y}\Bigr).
\end{equation}
The function \(\theta\) increases strictly and continuously from \(0\) to \((2m+1)\pi/2\), so
it takes each of the values \(\pi,2\pi,\dots,m\pi\) exactly once, at points
\(y_1<\dots<y_m\). These give \(m\) distinct roots \(\nu_k=-(1+y_k^2)/8<-\tfrac18\) of \(N\),
which are all its roots since \(\deg N=m\). By \eqref{eq:trig}, \(J(\nu_k)=(-1)^kR/2\). Hence
\(J\) changes sign on each interval \((\nu_{k+1},\nu_k)\), and since \(J(\nu_1)<0<J(0)\) by
\eqref{eq:leading}, also on \((\nu_1,0)\). This accounts for all \(m\) roots of \(J\) and
proves \eqref{eq:interlace}. The roots of \(J\) are simple, so \(N/J\) has the expansion
\eqref{eq:pf} with \(\sigma_k=N(\zeta_k)/J'(\zeta_k)\). Exactly \(k-1\) roots of \(N\) and
exactly \(k-1\) roots of \(J\) exceed \(\zeta_k\), and both leading coefficients are positive;
so \(N(\zeta_k)\) and \(J'(\zeta_k)\) both have sign \((-1)^{k-1}\), and \(\sigma_k>0\).
\end{proof}

For \(1\le k\le m\) let \(J_k(u)=J(u)/(u-\zeta_k)\). Since all \(\zeta_i\) are negative, \(J_k\)
is a polynomial of degree \(m-1\) with positive coefficients, and \(J_k(r)>0\) for \(r>0\).

\begin{theorem}[Coefficient positivity]\label{thm:positivity}
\textup{(i)} For every \(r>0\), all coefficients of \(Q\) and \(D_{\raw}\) in degrees
\(0,\dots,m\) are positive.

\textup{(ii)} Let \(\rho=\sqrt{1+8r}\). If \(\rho>\max_jx_j\) and
\(4r+\rho\ge\sum_jx_j\), then all coefficients of \(B_{\raw}\) in degrees \(0,\dots,m\) are
positive.
\end{theorem}
\begin{proof}
For \(m=0\) we have \(N=1\), \(J=d_0\), \(Q=d_0\), and \(B_{\raw}=2(1+d_0)(2r-d_0)\). Here
\(d_0=(x_0-1)/2\), and \(\rho>x_0\) gives \(2r=(\rho^2-1)/4>(x_0-1)(x_0+1)/4>d_0\). Let
\(m\ge1\).

(i) By \eqref{eq:pf}, \(uN/J=\alpha u+\sum_k\sigma_k+\sum_k\sigma_k\zeta_k/(u-\zeta_k)\). Hence
\[
 Q(u)=J(u)J_r\,\frac{\frac{uN(u)}{J(u)}-\frac{rN_r}{J_r}}{u-r}
 =J(u)J_r\Bigl(\alpha+\sum_k\frac{\sigma_k|\zeta_k|}{(u-\zeta_k)(r-\zeta_k)}\Bigr)
 =\alpha J_r\,J(u)+\sum_{k=1}^m\sigma_k|\zeta_k|\,J_k(r)\,J_k(u).
\]
Every term has nonnegative coefficients and the first has positive coefficients in all
degrees \(0,\dots,m\).

(ii) From \eqref{eq:pf}, \(N=\alpha J+\sum_k\sigma_kJ_k\). Substituting this and the formula
for \(Q\) into \(B_{\raw}=4rQ+4rN_rN-2J_r(N+J)\) and using \(J_r=(r+|\zeta_k|)J_k(r)\) gives
\begin{equation}\label{eq:Bdecomp}
 B_{\raw}=\alpha\Bigl[4rN_r-J_r\Bigl(\sum_jx_j+1-4r\Bigr)\Bigr]J
 +\sum_{k=1}^m\sigma_k\Bigl[4rN_r-2J_r\Bigl(1-\frac{2r|\zeta_k|}{r+|\zeta_k|}\Bigr)\Bigr]J_k ,
\end{equation}
where we used \(2+2/\alpha=\sum_jx_j+1\). Evaluating \eqref{eq:product} at \(x=\rho\), where
\(u=r\), gives \((\rho-1)N_r-2J_r=\Pi(\rho)>0\), that is,
\begin{equation}\label{eq:Jr}
 J_r<\tfrac12(\rho-1)N_r .
\end{equation}
Since \(4r=(\rho-1)(\rho+1)/2\), each bracket in the sum in \eqref{eq:Bdecomp} is at least
\(4rN_r-2J_r>\tfrac12(\rho-1)^2N_r>0\). If \(\sum_jx_j+1-4r\le0\), the first bracket is
positive as well. Otherwise, by \eqref{eq:Jr}, it exceeds
\[
 \tfrac12(\rho-1)N_r\Bigl[(\rho+1)-\Bigl(\sum_jx_j+1-4r\Bigr)\Bigr]
 =\tfrac12(\rho-1)N_r\Bigl(4r+\rho-\sum_jx_j\Bigr)\ \ge0 .
\]
Thus \(B_{\raw}\) is a combination of \(J,J_1,\dots,J_m\) with positive coefficients, and \(J\)
has positive coefficients in all degrees \(0,\dots,m\).
\end{proof}

Since the \(J_k\) have degree \(m-1\), the leading coefficient of \(B_{\raw}\) is
\(h_m\bigl[4rN_r-J_r(\sum_jx_j+1-4r)\bigr]\), by \eqref{eq:Bdecomp}. The proof therefore shows
that, for \(r>\max_ju_j\), the polynomial \(B_{\raw}\) has positive coefficients if and only
if its leading coefficient is positive. The second condition in \eqref{eq:rcondition} is a
convenient sufficient form of this, and it cannot simply be dropped: for the clustered roots
\(x=(1.059,\,1.098,\,1.003,\,1.029,\,1.038)\) one has \(\max_ju_j\approx0.0257\), but the
leading coefficient of \(B_{\raw}\) is negative at \(r=0.64\) and positive at \(r=0.66\), while
\eqref{eq:rcondition} holds from \(r\approx0.675\) on.

\begin{corollary}\label{cor:arith}
For \(x_j=j+2\) \((0\le j\le 2n-2)\), so that \(u_j=(j+1)(j+3)/8\), the polynomials
\(D_{\raw}\) and \(B_{\raw}\) have positive coefficients for every \(r>(4n^2-1)/8\).
\end{corollary}
\begin{proof}
Here \(\max_jx_j=2n\) and \(\sum_jx_j=(2n-1)(n+1)\). If \(\rho>2n\), then
\((\rho+1)^2>4n^2+4n+1>2\bigl(1+\sum_jx_j\bigr)=4n^2+2n\), and
\((\rho+1)^2\ge2(1+\sum_jx_j)\) is the same as \(4r+\rho\ge\sum_jx_j\).
\end{proof}

\section{Common totals and the count}\label{sec:assembly}
Fix \(r\) as in \eqref{eq:rcondition}; the two conditions there are those of
Theorem~\ref{thm:positivity}(ii), because \(r>u_j\) if and only if \(\rho>x_j\). Normalize and
choose the totals by
\begin{equation}\label{eq:normalize}
 D=\frac{D_{\raw}}{D_{\raw}(0)},\qquad B=\frac{B_{\raw}}{D_{\raw}(0)},\qquad
 A=(1+u)D,\qquad (\ET,\FT,\ST)=(2r,\,2,\,2r+2).
\end{equation}
Then \(A(0)=D(0)=1\), all coefficients of \(A,B,D\) in the ranges required by \eqref{eq:coeff}
are positive, and Proposition~\ref{prop:param}(ii) provides rate constants; with the
normalization \eqref{eq:rates} they are rational functions of the \(x_j\) and \(r\) with
rational coefficients. By Lemma~\ref{lem:identity}, for every \(j\),
\begin{equation}\label{eq:ratio}
 \frac{B(u_j)}{D(u_j)}=r+\frac{2(r-u_j)}{x_j-1}>r .
\end{equation}
For \(x=x_j\), \(u=u_j\) define
\begin{equation}\label{eq:sf}
 s=\frac{x-1}{2uD(u)},\qquad f=\frac{4}{x+1},\qquad\text{so that}\qquad E=uf=\frac{x-1}{2},
\end{equation}
and let \(z_j\) be the point \eqref{eq:states} with these \((u,s,f)\). It is a positive steady
state by Proposition~\ref{prop:param}(i). Its totals are computed from \eqref{eq:totalcoords}:
first \(suD=(x-1)/2\), so \(\FT=f(1+suD)=2\); next, by \eqref{eq:ratio},
\(suB=\tfrac12r(x-1)+(r-u)\), so
\[
 \ET=uf+f\bigl(\tfrac12r(x-1)+r-u\bigr)=\tfrac12 f\,r\,(x+1)=2r ;
\]
finally \(sA=(1+u)(x-1)/(2u)=4(1+u)/(x+1)=(1+u)f\), that is,
\begin{equation}\label{eq:freeinventory}
 \sum_{i=0}^nS_i=E+F ,
\end{equation}
and since the bound substrate is \(\ET+\FT-E-F\), we get \(\ST=\ET+\FT=2r+2\).

\begin{proof}[Proof of Theorem~\ref{thm:main} (count, prescribed ratios, rationality) and of
Theorem~\ref{thm:threshold}]
Given distinct \(u_j>0\), the numbers \(x_j=\sqrt{1+8u_j}\) are distinct and greater than
\(1\). Condition \eqref{eq:rcondition} holds for all sufficiently large \(r\). For any such
\(r\) the points \(z_0,\dots,z_{2n-2}\) are positive steady states of one and the same system
with the same totals, and they are distinct because their ratios \(E/F=u_j\) are. By
Theorem~\ref{thm:WS} there are no others. If all \(x_j\) are rational, take \(r\) rational;
then \(N,J,Q,B,D\), the rate constants \eqref{eq:rates}, and the concentrations
\eqref{eq:states} are rational. The statement about \(x_j=j+2\) is
Corollary~\ref{cor:arith}. Nondegeneracy is proved in Section~\ref{sec:nondeg}.
\end{proof}

\begin{remark}
(a) The map \(x\mapsto(x^2-1)/8\) is increasing on \((1,\infty)\), so lists of ratios with
rational \(x_j\) are dense among all lists; but we do not claim that an arbitrary list of
rational ratios \(u_j\) can be realized by rational data.
(b) Rescaling all concentrations by \(\kappa>0\) and the bimolecular constants
\(a_i,\alpha_i\) by \(1/\kappa\) gives the totals \(\kappa(2r,2,2r+2)\), and rescaling time
multiplies all rate constants by a common factor; neither changes the count. The symmetry
\(E\leftrightarrow F\), \(S_i\leftrightarrow S_{n-i}\) of \eqref{eq:network} gives examples
with \(\FT/\ET=r\) and prescribed ratios \(F/E\).
(c) The hypothesis \(r>\max_ju_j\) places all constructed steady states on the branch
\(u<r\), \(L>0\) of the chart \eqref{eq:chart}.
\end{remark}

\section{Nondegeneracy, a determinant formula, and stability}\label{sec:nondeg}

\subsection{Slopes}
\begin{proposition}\label{prop:slope}
In the construction of Section~\ref{sec:assembly}, for every \(j\),
\begin{equation}\label{eq:slope}
 H'(u_j)=-\frac{16\,(1+u_j)\,(\cs x_j+\es)}{(r-u_j)\,(x_j+1)^2\,D_{\raw}(u_j)}\;\Pi'(x_j),
 \qquad \cs x_j+\es=4rN_r+(x_j-1)J_r>0 .
\end{equation}
In particular \(H'(u_j)\ne0\), and if \(u_0<\dots<u_{2n-2}\) then
\(\sign H'(u_j)=(-1)^{j+1}\).
\end{proposition}
\begin{proof}
By \eqref{eq:ratio}, \(u_j\) lies in \(\mathcal U\cap\{u<r\}\), where
Lemma~\ref{lem:factor} applies. Dividing \eqref{eq:interpolation} by \((x-1)D_{\raw}(u)\)
gives
\[
 g(u):=\ell-\frac{2v}{x-1}=\frac{(\cs x+\es)\,\Pi(x)}{(x-1)\,D_{\raw}(u)},
\]
which vanishes at \(u_j\), and \(H-\ST=G\,g\) with
\(G=-2(1+u)\bigl(\ell+2v/(x+1)\bigr)/\bigl(\ell(\ell+v)\bigr)\). Hence
\(H'(u_j)=G(u_j)g'(u_j)\). At \(u_j\) we have \(\ell=2v/(x-1)\), so
\(\ell+2v/(x+1)=4vx/(x^2-1)\), \(\ell+v=v(x+1)/(x-1)\), and
\(G(u_j)=-4(1+u)x(x-1)/\bigl(v(x+1)^2\bigr)\). Since \(\Pi(x_j)=0\) and \(dx/du=4/x\),
\(g'(u_j)=4(\cs x+\es)\Pi'(x)/\bigl(x(x-1)D_{\raw}(u)\bigr)\). Multiplying gives
\eqref{eq:slope}. All factors other than \(\Pi'(x_j)=\prod_{k\ne j}(x_j-x_k)\) are positive,
and \(\Pi'(x_j)\) has sign \((-1)^{2m-j}=(-1)^j\) when the roots are ordered increasingly.
\end{proof}

\subsection{The restricted Jacobian}
We first fix coordinates. Put \(y=(S_1,\dots,S_n,C_1,\dots,C_n,Y_1,\dots,Y_n)\) and
\(z=(S_0,E,F)\), and order the totals as \((\ST,\ET,\FT)\). Then \(W=(W_y\;\;I_3)\) in the
coordinates \((y,z)\), so \(\ker W\) is the graph of \(z=-W_yy\) and \(y\) is a linear
coordinate system on it. The \(6n\) reaction vectors of \eqref{eq:network} lie in \(\ker W\)
and span it: in the coordinates \((C,Y,Z)\) on \(\ker W\), with \(Z\) as in \eqref{eq:pools},
a binding reaction changes only one complex coordinate, and the catalytic reactions at site
\(i\) change \(Z_i\) and no other \(Z_k\). So \(\ker W\) is the stoichiometric subspace. Let
\(\JS\) denote the matrix of \(DV|_{\ker W}\) in the coordinates \(y\), that is,
\(\JS=V_{y,y}-V_{y,z}W_y\), where \(V_y\) are the \(y\)-components of \(V\). A steady state is
nondegenerate if and only if \(\det\JS\ne0\).

\begin{proposition}[Determinant formula]\label{prop:det}
For arbitrary positive rate constants, at any positive steady state with \(u=E/F\ne r=\ET/\FT\),
\begin{equation}\label{eq:det}
 \det\JS=(-1)^{n+1}\Bigl(\prod_{i=1}^n(b_i+c_i)\,\alpha_i\,\gamma_i\Bigr)\,F^{\,n}\,u\,M(u)\,H'(u),
\end{equation}
where \(H\) is formed with the enzyme totals of that steady state. In particular the steady
state is nondegenerate if and only if \(H'(u)\ne0\).
\end{proposition}
\begin{proof}
By the Schur complement with respect to the block \(I_3\),
\[
 \det\JS=\det\begin{pmatrix}V_{y,y}&V_{y,z}\\ W_y&I_3\end{pmatrix}.
\]
Let \(g=(g_Z,g_C,g_Y)\) consist of the \(3n\) functions \(g_{Z,i}=c_iC_i-\gamma_iY_i\),
\(g_{C,i}=\dot C_i\), \(g_{Y,i}=\dot Y_i\). By \eqref{eq:pools}--\eqref{eq:poolcurrent},
\(g=UV_y\) for a constant unitriangular matrix \(U\), since \(g_{Z,i}\) is \(\dot S_i\) plus
components of \(V_y\) that come later in the ordering of \(y\). Hence the first \(3n\) rows can
be replaced by \((g_y\;\;g_z)\) without changing the determinant. In the block form belonging
to \(y=(S;C,Y)\),
\[
 g_y=\begin{pmatrix}0&\ast\\ \ast&\Delta\end{pmatrix},\qquad
 \Delta=\operatorname{diag}\bigl(-(b_i+c_i);\,-(\beta_i+\gamma_i)\bigr),
\]
and the Schur complement of \(\Delta\) is the Jacobian with respect to \(S_1,\dots,S_n\) of
the reduced currents \(c_ip_iS_{i-1}E-\gamma_iq_iS_iF\). It is lower bidiagonal with diagonal
entries \(-\gamma_iq_iF\). Therefore
\[
 \det g_y=\prod_i(b_i+c_i)(\beta_i+\gamma_i)\cdot(-1)^n\prod_i\gamma_iq_iF
 =(-1)^nF^n\prod_i(b_i+c_i)\alpha_i\gamma_i\ne0 .
\]
In particular the zero set of \(g\) is locally the graph of a function \(y=\varphi(z)\), in
agreement with Proposition~\ref{prop:param}, and
\[
 \det\begin{pmatrix}g_y&g_z\\ W_y&I_3\end{pmatrix}
 =\det g_y\cdot\det\bigl(I_3-W_yg_y^{-1}g_z\bigr)
 =\det g_y\cdot\det\frac{\partial(\ST,\ET,\FT)}{\partial(S_0,E,F)},
\]
where the last matrix is the derivative of the totals along the steady-state manifold
\(z\mapsto(\varphi(z),z)\), because \(\varphi_z=-g_y^{-1}g_z\). We compute it in the
coordinates \((u,s,f)=(E/F,S_0,F)\), for which
\(\det\partial(u,s,f)/\partial(S_0,E,F)=-1/F\). From \eqref{eq:totalcoords},
\begin{equation}\label{eq:pivot}
 \det\frac{\partial(\ET,\FT)}{\partial(s,f)}
 =\det\begin{pmatrix}ufB&u(1+sB)\\ fuD&1+suD\end{pmatrix}=uf\,(B-uD)=uf\,M(u)\ne0,
\end{equation}
so near the given steady state the enzyme totals determine \((s,f)\) as functions of \(u\),
namely \eqref{eq:chart}, and the Schur complement of the block \eqref{eq:pivot} in
\(\partial(\ST,\ET,\FT)/\partial(u,s,f)\) is the derivative of \(\ST\) along that curve, which
is \(H'(u)\). Hence \(\det\partial(\ST,\ET,\FT)/\partial(u,s,f)=ufM(u)H'(u)\), and collecting
the factors gives \eqref{eq:det}.
\end{proof}

\begin{proof}[Proof of Theorem~\ref{thm:main} (nondegeneracy)]
Combine Propositions~\ref{prop:slope} and~\ref{prop:det}.
\end{proof}

\begin{proof}[Proof of Corollary~\ref{cor:open}]
Let \(k^\circ\) and \(T^\circ\) be the rate constants and totals of Theorem~\ref{thm:main}. In
the coordinates \(y\) on the compatibility classes, the steady-state equations are
\(V_y(y,T-W_yy;k)=0\), a system of \(3n\) polynomial equations in \(3n\) unknowns depending
smoothly on \((k,T)\in\R^{6n+3}_{>0}\), and its derivative in \(y\) is \(\JS\). By the
implicit function theorem each \(z_j\) continues to a smooth family of nondegenerate positive
steady states for \((k,T)\) in a neighborhood \(\mathcal O_j\) of \((k^\circ,T^\circ)\), and
after shrinking the \(\mathcal O_j\) the continued states remain pairwise distinct. On
\(\bigcap_j\mathcal O_j\) there are at least \(2n-1\) positive steady states, all
nondegenerate, and by Theorem~\ref{thm:WS} there are no more. The perturbed parameters need
not satisfy \(A=(1+u)D\) or \(\ST=\ET+\FT\).
\end{proof}

\begin{corollary}\label{cor:unstablegeneral}
For arbitrary positive rate constants, a positive steady state with \(u\ne r\) and
\((r-u)H'(u)>0\) is unstable; more precisely \(\JS\) has an odd number of positive real
eigenvalues. If \((r-u)H'(u)<0\), this number is even.
\end{corollary}
\begin{proof}
When \(\det\JS\ne0\), its sign is \((-1)^{e_-}\), where \(e_-\) is the number of negative
real eigenvalues counted with multiplicity, because nonreal eigenvalues come in conjugate
pairs. For the same reason \(3n-e_-\) has the parity of the number \(e_+\) of positive real
eigenvalues. So \(\sign\det\JS=(-1)^{n}(-1)^{e_+}\). By
\eqref{eq:det} and since \(M(u)\) has the sign of \(L(u)\), hence of \(r-u\), we also have
\(\sign\det\JS=(-1)^{n+1}\sign\bigl((r-u)H'(u)\bigr)\).
\end{proof}

\begin{proof}[Proof of Theorem~\ref{thm:unstable}]
All constructed steady states have \(u_j<r\), and \(\sign H'(u_j)=(-1)^{j+1}\) by
Proposition~\ref{prop:slope}. Apply Corollary~\ref{cor:unstablegeneral}.
\end{proof}

Corollary~\ref{cor:unstablegeneral} is the form taken, for this network, by the general
principle that the signs of \(\det(-\JS)\) at the nondegenerate steady states of a dissipative
network without boundary steady states add up to one \cite{CFMW2017}; here it is obtained by
direct computation, together with the exact value of the determinant. It does not decide the
stability of the even-indexed steady states, which depends on the \(3n\) kinetic parameters
that do not enter the steady-state equations.

\section{Further properties of the constructed steady states}\label{sec:properties}
The formulas of Section~\ref{sec:assembly} give closed expressions for several quantities of
biochemical interest. We use the normalization \eqref{eq:rates} and order the ratios
increasingly.

\begin{proposition}\label{prop:properties}
At the constructed steady state with auxiliary coordinate \(x=x_j\):
\begin{enumerate}
\item[\textup{(i)}] \(E=(x-1)/2\), \(F=4/(x+1)\), and the total free substrate equals \(E+F\).
The fractions of kinase and phosphatase bound in complexes are \(1-(x-1)/(4r)\) and
\(1-2/(x+1)\).
\item[\textup{(ii)}] The fully phosphorylated substrate is
\(S_n=\tfrac12D_{n-1}\,(x-1)\,u^{n-1}/D(u)\), and \(S_n+Y_n=S_n\,(x+5)/(x+1)\). Both are
strictly increasing functions of \(j\).
\item[\textup{(iii)}] The total phosphorylation flux is
\(\sum_ic_iC_i=\sum_i\gamma_iY_i=2(x-1)/(x+1)\), which is strictly increasing in \(j\).
\end{enumerate}
\end{proposition}
\begin{proof}
(i) is \eqref{eq:sf}, \eqref{eq:freeinventory}, and \(\sum C_i=\ET-E\), \(\sum Y_i=\FT-F\).
(ii) Since \(A=(1+u)D\), \(t_n=D_{n-1}\), so \(S_n=t_nu^ns\) and \eqref{eq:sf} give the first
formula, and \(Y_n=D_{n-1}u^nsf=fS_n\) gives the second. The function
\(D(u)/u^{n-1}=\sum_iD_iu^{i-(n-1)}\) is nonincreasing, \(x-1\) is increasing, and
\((x-1)(x+5)/(x+1)=x+3-8/(x+1)\) is increasing. (iii) With \(\gamma_i=1\) the flux is
\(\sum Y_i=2-F\).
\end{proof}

So the \(2n-1\) steady states are distinguished by an output that is commonly measured, the
amount of fully phosphorylated substrate, and they are ordered in the same way by that
output, by the free-enzyme ratio, and by the rate at which phosphate is turned over. The
steady states are not thermodynamic equilibria: opposing catalytic currents are positive and
equal, and in a model with explicit ATP the flux in (iii) would be the rate of ATP
consumption needed to hold the state. By (i), in these examples the amount of enzyme equals
the amount of substrate and most of the kinase is bound; they are far from the
substrate-excess regime in which Michaelis--Menten reductions apply.

\begin{remark}[Large \(r\)]\label{rem:larger}
For fixed \(x_j\) and \(r\to\infty\), the coefficient formulas of Appendix~\ref{app:bound}
give \(D\to J/d_0\) and \(B/r\to(J+N)/d_0\) coefficientwise. Hence \(a_i\) and \(\alpha_i\)
converge to positive limits, while \(b_i=c_i\) is asymptotic to
\(r^{-1}d_{i-1}/(d_{i-1}+h_{i-1})\). Increasing the kinase excess therefore slows the kinase
catalytic step relative to the phosphatase in these examples, and the free concentrations and
the slopes \eqref{eq:slope} have finite nonzero limits.
\end{remark}

\section{Examples and exact stability computations}\label{sec:examples}

\subsection{A three-site example}
Let \(n=3\) and \((x_0,\dots,x_4)=(2,3,4,5,6)\), so that the prescribed ratios are
\(u_j=\tfrac38,1,\tfrac{15}8,3,\tfrac{35}8\). The recurrence gives
\[
 N=1200+1256u+64u^2,\qquad J=60+1852u+608u^2 .
\]
By Corollary~\ref{cor:arith} every \(r>35/8\) is admissible. For \(r=5\) we obtain
\begin{equation}\label{eq:example}
 D=1+\frac{69002}{3405}u+\frac{9808}{3405}u^2,\qquad
 B=\frac{34797}{454}+\frac{308833}{2270}u+\frac{5236}{1135}u^2,\qquad A=(1+u)D,
\end{equation}
the totals \((\ET,\FT,\ST)=(10,2,12)\), and the rate constants of Table~\ref{tab:rates}. The
five steady states are listed in Table~\ref{tab:states}; all twelve concentrations of each are
rational numbers given by \eqref{eq:states} and \eqref{eq:sf}, for instance
\(E=\tfrac12,1,\tfrac32,2,\tfrac52\) and \(F=\tfrac43,1,\tfrac45,\tfrac23,\tfrac47\).
Figure~\ref{fig:residual}(a) shows the function \(H-\ST\).

\begin{table}[ht]
\centering
\caption{Rate constants of the three-site example \((r=5)\).}\label{tab:rates}
\begin{tabular}{ccccc}
\toprule
Site \(i\) & \(a_i\) & \(b_i=c_i\) & \(\alpha_i\) & \(\beta_i=\gamma_i\)\\
\midrule
1 & \(2\) & \(454/34797\) & \(6810/72407\) & \(1\)\\
2 & \(138004/72407\) & \(138004/926499\) & \(69002/39405\) & \(1\)\\
3 & \(9808/39405\) & \(2452/3927\) & \(2\) & \(1\)\\
\bottomrule
\end{tabular}
\end{table}

\begin{table}[ht]
\centering
\caption{The five positive steady states of the three-site example. Decimal entries are
rounded values of exact rational numbers. The last column is the number of eigenvalues of
\(\JS\) in the open right half-plane, obtained from an exact Routh array; no eigenvalue lies
on the imaginary axis.}\label{tab:states}
\begin{tabular}{cccccccc}
\toprule
\(j\) & \(u=E/F\) & \(E\) & \(F\) & \(S_3\) & \(S_3+Y_3\) & \(H'(u_j)\) & unstable\\
\midrule
0 & \(3/8\) & \(1/2\) & \(4/3\) & \(0.02249\) & \(0.05248\) & \(-0.13324\) & 0\\
1 & \(1\) & \(1\) & \(1\) & \(0.11930\) & \(0.23859\) & \(\phantom{-}0.01315\) & 1\\
2 & \(15/8\) & \(3/2\) & \(4/5\) & \(0.30922\) & \(0.55660\) & \(-0.00561\) & 0\\
3 & \(3\) & \(2\) & \(2/3\) & \(0.59107\) & \(0.98512\) & \(\phantom{-}0.00780\) & 1\\
4 & \(35/8\) & \(5/2\) & \(4/7\) & \(0.95195\) & \(1.49591\) & \(-0.06498\) & 0\\
\bottomrule
\end{tabular}
\end{table}

\begin{figure}[ht]
\centering
\includegraphics[width=\linewidth]{figures/residuals.pdf}
\caption{The scalar function \(H(u)-\ST\) on the branch \(u<r\) for the family \(x_j=j+2\):
(a) \(n=3\), \(r=5\); (b) \(n=5\), \(r=13\). The vertical axis is
\(\operatorname{arcsinh}\) of the rescaled residual, labelled by the residual itself. Filled
markers are asymptotically stable steady states, open markers are saddles
(Proposition~\ref{prop:census}). The curves are drawn from exact rational evaluations; the
count is certified by the formulas, not by the plot.}\label{fig:residual}
\end{figure}

Nondegeneracy guarantees that the count persists under small perturbations but gives no size
for the neighborhood. The following statement, for one parameter direction, is an example of
what can be certified by exact evaluation.

\begin{proposition}\label{prop:window}
With the rate constants of Table~\ref{tab:rates} and \((\ET,\FT)=(10,2)\), the system has
exactly five positive steady states for every \(\ST\in[12-\tfrac1{500},\,12+\tfrac1{500}]\).
\end{proposition}
\begin{proof}
Here \(L=B-5D\) has a negative leading coefficient and \(L(0),L(5)>0\), so \(L>0\) on
\([0,5]\) by concavity, and \(H\) is continuous on \((0,5)\). Exact evaluation at
\(u=\tfrac1{20},\tfrac{69}{128},\tfrac{4327}{3200},\tfrac{7917}{3200},\tfrac{6373}{1600},
\tfrac{399}{80}\) gives values of \(H-12\) that alternate in sign, starting with a positive
one, and all exceed \(\tfrac1{500}\) in absolute value (approximately
\(1.02,\,-0.0080,\,0.0021,\,-0.0021,\,0.0092,\,-2.6\)). The intermediate value theorem gives five
solutions of \(H(u)=\ST\) in \((0,5)\), hence five positive steady states, and
Theorem~\ref{thm:WS} excludes a sixth.
\end{proof}

The window is narrow, about \(0.03\%\) of \(\ST\), because the three interior extrema of
\(H-\ST\) are small. Correspondingly the relative sensitivities
\(|d\log u/d\log\ST|=\ST/(u|H'(u)|)\) at the five steady states are approximately
\(240,\,913,\,1140,\,513,\,42\), and the leading eigenvalues of \(\JS\), computed in floating
point, are of order \(10^{-4}\) (approximately \(-5.8\times10^{-4}\), \(8.8\times10^{-5}\),
\(-4.5\times10^{-5}\), \(5.8\times10^{-5}\), \(-1.9\times10^{-4}\)) while the remaining
eigenvalues are of order \(10^{-1}\) to \(1\). The example is a certified benchmark with exact
data; it is not tuned for wide separation of the steady states, and we have not tried to
optimize the \(x_j\), \(r\), or the \(3n\) free kinetic parameters for that purpose.

\subsection{Stability for \texorpdfstring{\(n\le10\)}{n at most 10}}
\begin{proposition}[Computer-assisted]\label{prop:census}
Let \(1\le n\le10\), \(x_j=j+2\), \(r=\lfloor(4n^2-1)/8\rfloor+1\), and let the rate constants
be given by \eqref{eq:normalize} and \eqref{eq:rates}. Then the \(n\) steady states
\(z_0,z_2,\dots,z_{2n-2}\) are asymptotically stable within their compatibility class: all
\(3n\) eigenvalues of \(\JS\) have negative real part. At each of the \(n-1\) steady states
\(z_1,z_3,\dots,z_{2n-3}\), exactly one eigenvalue of \(\JS\) has positive real part and none
lies on the imaginary axis.
\end{proposition}
\begin{proof}
For each of the \(100\) steady states the matrix \(\JS\) is rational. Its characteristic
polynomial was computed exactly by the Faddeev--LeVerrier recursion in integer arithmetic,
with the Cayley--Hamilton identity as a check, and the Routh array \cite[Ch.~XV]{Gantmacher}
was formed in rational arithmetic. In every case all \(3n+1\) entries of the first column are
nonzero, so the array is regular, there are no eigenvalues on the imaginary axis, and the
number of sign changes in the first column is the number of eigenvalues in the open right
half-plane. It is \(0\) for even \(j\) and \(1\) for odd \(j\).
\end{proof}

For \(n=3\) and \(n=4\), examples with \(2n-1\) steady states were first given in
\cite{FHC2014}. The numbers of right half-plane eigenvalues in Proposition~\ref{prop:census}
agree with the parities required by Theorem~\ref{thm:unstable}.

\begin{conjecture}\label{conj:stable}
For every \(n\ge1\) there are rate constants and totals for which \eqref{eq:network} has
\(2n-1\) positive steady states of which \(n\) are asymptotically stable. More specifically,
the conclusion of Proposition~\ref{prop:census} holds for all \(n\).
\end{conjecture}

\section{Formal verification and reproducibility}\label{sec:formal}
The lower bound has been formalized in Lean~4 \cite{Lean4} using Mathlib \cite{Mathlib}. The
development defines the state space, the six rate constants per site, the mass-action vector
field of \eqref{eq:network} componentwise, and the three totals, and proves
\begin{quote}
\texttt{theorem PhosphorylationSharpness.sharp\_lower\_bound}\\
\texttt{\ \ \ \ (n :\ $\mathbb N$) (hn :\ 1 $\le$ n) :\ Attains n (2*n-1)}
\end{quote}
where \texttt{Attains n k} states that there exist positive rate constants, positive totals
\(\ET,\FT,\ST\), and an injective family of \(k\) states, each with all coordinates positive,
at which every component of the vector field vanishes and the three totals take the given
values. The statement thus concerns the full system of \(3n+3\) equations and contains no
reduced equation, root count, or numerical input. The proof follows
Sections~\ref{sec:model}--\ref{sec:assembly} for the family \(x_j=j+2\), with one difference:
positivity of the coefficients of \(D_{\raw}\) and \(B_{\raw}\) is proved for all sufficiently
large \(r\) by the elementary coefficient estimate of Appendix~\ref{app:bound}, in the form of
a limit as \(r\to\infty\), and not by Theorem~\ref{thm:positivity}. It uses raw rather than
normalized coefficients; multiplying \(A,B,D\) by a positive constant and dividing \(s\) by it
changes neither the rate constants \eqref{eq:ratesgeneral} nor the concentrations
\eqref{eq:states}. The five modules (\texttt{Recurrence}, \texttt{Conversion},
\texttt{Source}, \texttt{Assembly}, \texttt{Resolution}; about 580 lines) compile without
warnings under Lean \texttt{v4.30.0} with warnings treated as errors, and the theorem depends
only on the axioms \texttt{propext}, \texttt{Classical.choice}, and \texttt{Quot.sound}. A
further module, \texttt{PublicationAlgebra}, verifies the recurrence identity
\eqref{eq:product} and coefficient positivity for arbitrary lists of roots greater than one,
the determinant identity \eqref{eq:pivot}, the factorization \eqref{eq:residual} of the
numerator, and the sign of the prefactor in \eqref{eq:slope}.

The following parts of the paper are proved only in the conventional way and are not covered
by the formal development: the upper bound (Theorem~\ref{thm:WS}, quoted from
\cite{WangSontag2008}), arbitrary prescribed ratios and rationality in
Theorem~\ref{thm:main}, the interlacing lemma and Theorem~\ref{thm:positivity}, the slope and
determinant formulas, Corollary~\ref{cor:open}, and Theorem~\ref{thm:unstable}.
Propositions~\ref{prop:window} and~\ref{prop:census} are exact computations that are not
formalized.

{\sloppy The source package of this paper contains the Lean modules and a Python script,
\texttt{check\_paper.py}, which uses only the standard library and rational arithmetic. It
reconstructs the examples for \(n\le8\) for the family \(x_j=j+2\) and checks, for all \(64\)
steady states, that the vector field vanishes, that the totals agree, and that the slope
formula \eqref{eq:slope} holds; it checks Lemma~\ref{lem:interlace} by Sturm sequences and
Theorem~\ref{thm:positivity} at the boundary of the admissible range on several hundred
random root configurations; it checks the determinant formula \eqref{eq:det} exactly for the
constructed steady states with \(n\le4\) and at random steady states of random rate
constants; and it reproduces every number in Section~\ref{sec:examples}, including
Proposition~\ref{prop:census} (for \(n\le6\) by default, and for \(n\le10\) when run with
the option \texttt{-{}-full}). These computations test the exposition and the implementation. The proofs
above do not depend on them, with the exception of Propositions~\ref{prop:window}
and~\ref{prop:census}.\par}

\section{Discussion and open problems}\label{sec:discussion}
The maximal number of positive steady states of the sequential distributive \(n\)-site
phosphorylation cycle is \(2n-1\) for every \(n\), it is attained on an open set of
parameters, and the locations of the steady states on the axis of the free-enzyme ratio can
be prescribed arbitrarily. The proof is constructive and produces rational examples of any
size in \(O(n^2)\) arithmetic operations.

The construction has clear limitations, and we list the questions they suggest.

\emph{Stability.} We prove that \(n-1\) of the constructed steady states are unstable and
verify for \(n\le10\) that the other \(n\) are stable. A proof of
Conjecture~\ref{conj:stable} would require control of all \(3n\) eigenvalues. The kinetic
parameters \(b_i,\beta_i,\gamma_i\) of Proposition~\ref{prop:param}(ii) do not affect the
steady states and may be chosen to help, for instance in a limit of fast binding, where the
dynamics reduces to the \(n\) slow variables \(Z_i\) of \eqref{eq:pools}. The argument of
\cite{FRW2020} relies on a regime that cannot contain more than \(2\lfloor n/2\rfloor+1\)
steady states, so a different reduction is needed.

\emph{Enzyme levels.} In our examples \(\ST=\ET+\FT\), and the ratio \(\ET/\FT=r\) allowed by
\eqref{eq:rcondition} is large: of order \(n^2\) for equally spaced \(x_j\), and about
\(n/2\) when the \(x_j\) cluster near \(1\). Wang and Sontag showed that, generically, there are at
most \(n+1\) steady states when the enzymes are scarce relative to the substrate, and only
one when they are in large excess \cite[Theorems~4 and~5]{WangSontag2008}. It would be interesting to know how small the ratios \(\ET/\ST\)
and \(\FT/\ST\) can be in a system with \(2n-1\) steady states, and whether the enzyme totals
can be comparable to each other.

\emph{Size of the region.} Corollary~\ref{cor:open} gives an open set without an estimate of
its size, and Proposition~\ref{prop:window} shows that in the example it is small in at least
one direction. The descriptions of multistationarity regions in
\cite{CFMW2017,BihanDickensteinGiaroli2020,GRPD2019,FKdWY2022} are explicit but concern fewer
steady states. A quantitative description of the region with \(2n-1\) steady states, even
for \(n=3\), is open. The freedom in the choice of the \(x_j\) and of \(r\) can presumably be
used to enlarge it.

\emph{Other mechanisms.} The method uses two features of \eqref{eq:network}: the monomial
parametrization of the steady states, and a parameter locus on which the scalar steady-state
equation factors after a quadratic change of variable. It is natural to ask whether a
similar locus exists for related networks, such as cycles with several kinases or
phosphatases, mixed processive--distributive mechanisms, or cascades, for which sharp bounds
are not known.

\appendix
\section{An elementary positivity bound}\label{app:bound}
The following estimate gives the positivity of \(D_{\raw}\) and \(B_{\raw}\) for large \(r\)
without the interlacing lemma. It is the route taken in the formal proof. Let
\(A_\Sigma=\sum_ih_i\), \(D_\Sigma=\sum_id_i\), \(d_\ast=\min_id_i\), and \(\eta=h_m\).

\begin{lemma}\label{lem:crude}
If
\begin{equation}\label{eq:rbound}
 r>\max\Bigl\{1,\ \frac{4A_\Sigma D_\Sigma}{\eta d_\ast},\
 \frac{2D_\Sigma(A_\Sigma+D_\Sigma)}{\eta d_\ast}\Bigr\},
\end{equation}
then \(Q\), \(D_{\raw}\), \(B_{\raw}\) have positive coefficients in degrees \(0,\dots,m\).
\end{lemma}
\begin{proof}
Write \(Q=\sum_iq_iu^i\). Comparing coefficients in \((u-r)Q=uJ_rN-rN_rJ\) gives
\(q_0=N_rd_0\) and \(q_i=N_rd_i+(q_{i-1}-J_rh_{i-1})/r\), hence
\begin{equation}\label{eq:qformula}
 q_i=N_rd_i+\sum_{k<i}r^{k-i}\bigl(N_rd_k-J_rh_k\bigr).
\end{equation}
Since \(r>1\), \(\eta r^m\le N_r\le A_\Sigma r^m\) and \(J_r\le D_\Sigma r^m\), and
\(r^{k-i}\le r^{-1}\) for \(k<i\). So the sum in \eqref{eq:qformula} is at most
\(r^{-1}(N_rD_\Sigma+J_rA_\Sigma)\le2A_\Sigma D_\Sigma r^{m-1}\) in absolute value, and the
first bound in \eqref{eq:rbound} gives \(q_i>\tfrac12\eta d_\ast r^m\). The coefficient of
\(u^i\) in \(B_{\raw}\) is
\(4rq_i+4rN_rh_i-2J_r(h_i+d_i)>2\eta d_\ast r^{m+1}-2D_\Sigma(A_\Sigma+D_\Sigma)r^m>0\) by
the second bound.
\end{proof}

The individual terms in \eqref{eq:qformula} can be negative, so positivity of \(Q\) is not
visible term by term; Theorem~\ref{thm:positivity}(i) shows that the sum is nevertheless
positive for every \(r>0\). For the three-site example the right side of \eqref{eq:rbound} is
\(6615\), whereas Theorem~\ref{thm:threshold} admits every \(r>35/8\). Formula
\eqref{eq:qformula} also gives the limits used in Remark~\ref{rem:larger}:
\(q_i/r^m\to\eta d_i\) and \((B_{\raw})_i/r^{m+1}\to4\eta(d_i+h_i)\), while
\(D_{\raw}(0)=4N_rd_0\).

\bibliographystyle{abbrv}
\bibliography{refs}
\end{document}
